% Bagean: sajarah
% Bab: biografi
% Pérangan: ernst-zermelo

\documentclass[../../../include/open-logic-section]{subfiles}

\begin{document}

\olfileid{his}{bio}{zer}

\olsection{Ernst Zermelo}

\olphoto{zermelo-ernst}{Ernst Zermelo}

Ernst Zermelo lair tanggal 27 Juli 1871
ing Berlin, Jerman. Dhèwèké nduwèni
limang sadulur wadon, nanging kulawargané
kerep nandhang lara lan mung telu kang
urip nganti diwasa. Wong tuwané uga
séda nalika dhèwèké isih enom,
saéngga dhèwèké lan para saduluré
dadi bocah yatim piatu nalika umur
pitung welas taun. Zermelo kasengsem
banget marang seni, mligi geguritan.
Dhèwèké kawentar landhep pamikiré,
pinter guyon, lan kritis. Prestasi
matématikané kang paling misuwur
kalebu ngenalaké aksioma pilihan
(ing taun 1904) lan aksiomatisasi
téori himpunan (ing taun 1908).

Minaté Zermelo ing universitas
warna-warna. Dhèwèké njupuk
kuliah fisika, matématika, lan
filsafat. Ing sangisoré bimbingan
Hermann Schwarz, Zermelo
ngrampungaké disertasiné
\emph{Panlitèn ing Kalkulus Variasi}
ing taun 1894 ing Universitas
Berlin. Ing taun 1897,
dhèwèké mutusaké nerusaké
studi ing Universitas
G\"{o}ttingen, ing kono
dhèwèké kapengaruh banget
déning karya dhasar
David Hilbert. Ing taun
1899 dhèwèké wis
nyukupi syarat dadi
profèsor, nanging
durung olèh jabatan
iku nganti sewelas
taun sabanjuré---
mbokmenawa amarga
polatané kang ora
lumrah lan ``ketergesan
amarga gugup.''

Zermelo pungkasané
olèh jabatan profèsor
kanthi bayaran ing
Universitas Zurich
ing taun 1910,
nanging kepeksa
pensiun ing taun
1916 amarga
tuberkulosis.
Sawisé mari,
dhèwèké diwènèhi
jabatan profèsor
pakurmatan ing
Universitas Freiburg
ing taun 1921.
Ing mangsa iki
dhèwèké nggarap
dhasar-dhasar
matématika. Dhèwèké
ora sarujuk banget
karo karya Thoralf
Skolem lan Kurt
G\"{o}del, lan
ngritik cara-carané
kanthi terang-terangan
ing makalah-makalahé.
Dhèwèké dipecat
saka jabatané ing
Freiburg ing taun
1935, amarga ora
disenengi lan
amarga nentang
munggahé Hitler
marang kuwasa
ing Jerman.

Taun-taun pungkasan
uripé Zermelo
diwarnai déning
kasepian. Sawisé
dipecat ing taun
1935, dhèwèké
ninggalaké matématika.
Dhèwèké pindhah
menyang désa lan
urip prasaja.
Dhèwèké omah-omah
ing taun 1944,
lan saya gumantung
marang garwané
amarga paningalé
saya ilang.
Zermelo wuta
babar pisan
ing taun 1951.
Dhèwèké séda
ing G\"{u}nterstal,
Jerman, tanggal
21 Mei 1953.

\begin{reading}
Kanggo biografi
Zermelo kang
jangkep, delengen
\citet{Ebbinghaus2015}.
Makalah-makalah
utama Zermelo
taun 1904 lan
1908 bisa diwaca
ing basa Jerman
asliné
\citep{Zermelo1904,Zermelo1908}.
Karya-karya
Zermelo kang
diklumpukaké,
kalebu tulisané
ngenani fisika,
kasedhiya ing
terjemahan
basa Inggris
ing
\citep{Ebbinghaus2010,Ebbinghaus2013}.
\end{reading}

\end{document}
